\documentclass[12pt,a4paper]{article}
\usepackage[utf8]{inputenc}
\usepackage[T2A]{fontenc}
\usepackage[russian]{babel}
\usepackage{graphicx}
\usepackage{booktabs}
\usepackage{longtable}
\usepackage{geometry}

\usepackage{float}

\geometry{top=2cm, left=2cm, right=2cm, bottom=2cm}
\usepackage{paracol}
\usepackage{src/title}
\usepackage{hyperref}

\usepackage{listings}
\usepackage{xcolor}

\title{%TODO название лабы}
\author{Ануфриев А.С. 465029}
\supervisor{
%TODO практик
}
\date{10.10.2026}

\abstract{
%TODO кратное описание
%}
\keywords{%TODO ключевые слова}

\begin{document}
\maketitle
%TODO далее содержание согласно заданию, например:

\chapter*{Шаблон отчёта по Модулю 1}

\begin{enumerate}

\item Введение:
\begin{itemize}

\item Описание целей проекта и его значимости.

\end{itemize}



\item Анализ требований:
\begin{itemize}

\item Определение основных требований к системе поддержки принятия решений.

\item Выявление требований к базе знаний и онтологии для представления знаний.

\end{itemize}



\item Изучение основных концепций и инструментов:
\begin{itemize}

\item Обзор основных концепций баз знаний и онтологий.

\item Изучение Prolog и его возможностей для разработки систем искусственного интеллекта.

\item Ознакомление с инструментами и библиотеками, подходящими для работы с базами знаний и онтологиями на Prolog.

\end{itemize}



\item Реализация системы искусственного интеллекта (системы поддержки принятия решений):
\begin{itemize}

\item Создание правил и логики вывода для принятия решений на основе базы знаний и онтологии.

\item Тестирование и отладка системы, обеспечение ее функциональности и эффективности.

\end{itemize}



\item Оценка и интерпретация результатов:
\begin{itemize}

\item Примеры запросов для БЗ и онтологии, сравнение разницы реализации.

\item Оценка соответствия системы поставленным требованиям и достижению целей проекта.

\item Интерпретация результатов и описание дальнейших возможностей развития и улучшения системы.

\end{itemize}



\item Заключение:
\begin{itemize}

\item Описание преимуществ и потенциальных применений разработанной системы искусственного интеллекта на базе Prolog, баз знаний и онтологий.

\end{itemize}



\end{enumerate}



\end{document}